%%%%%%%%%%%%%%%%%%%%%%%%%%%%%%%%%%%%%%%%%%%%%%%%%%%%%%%%%%%%
%% CHAPTER 9 -- Teaching Manners and the Cost of Doing So
%%%%%%%%%%%%%%%%%%%%%%%%%%%%%%%%%%%%%%%%%%%%%%%%%%%%%%%%%%%%

\chapter{Teaching Manners and the Cost of Doing So}
\label{ntg:ch9}

\scenelabel
\begin{scenestrip}
\sceneitem{The problem}
\sceneitem{How we got here}
\sceneitem{Where we stand}
\sceneitem{What goes wrong}
\end{scenestrip}

\paragraph{The problem.} Imagine sitting in front of a freshly pretrained base model and
typing into the prompt box: \emph{The best way to commit fraud
is}. The model, eager to do its job, will continue the
sentence. It will produce coherent prose. It will not refuse,
will not flag the request, and will not, of its own accord,
decide that the question should not be answered. This is not
the model misbehaving. The model is doing exactly what
Chapter~6 trained it to do, which is to produce the most
plausible continuation of the prefix it was given.

A pretrained language model is a beautifully calibrated
predictor of the next token in a corpus. That is what we
asked it to be, and it is not the same thing as a helpful
assistant. A base model will gladly continue an essay that
begins ``the best way to commit fraud is''. It will produce
a confident answer to a question whose answer it does not
know. It will agree with the most recent claim in the
conversation regardless of whether the claim is true. It will
generate a recipe for any plausible-sounding dish, including
ones that would poison the cook. None of these are bugs in
the pretraining objective. They are what the pretraining
objective produces.

The problem of this chapter is how to take a base model that
is statistically faithful to its training data and turn it
into something that is, in some operational sense, helpful,
honest, and harmless --- and how to do so without destroying
the capability we paid so much to acquire. The answer is a
discipline called \emph{alignment}, and it is the part of the
modern pipeline whose successes and failures are most visible
to end users.

\paragraph{How we got here.} The first instinct, when researchers began trying to make
pretrained models behave usefully, was the obvious one. They
fine-tuned the model on examples of the desired behaviour.
Show the model a thousand prompt-response pairs in which the
response is the kind a good assistant would give. Train the
model to produce responses like the example responses. The
technical name is \emph{supervised fine-tuning}, abbreviated
SFT, and it is still the first stage of essentially every
modern post-training pipeline.

SFT works for some kinds of behaviour. It works for teaching
the model the format of an instruction-response interaction.
It works for teaching the model to refuse a particular class
of request with particular wording. It is brittle for
nuanced preferences. It is hard to write a single ``correct''
answer to a question like ``how should I phrase this difficult
email to my boss?''. The space of acceptable answers is wide;
the space of slightly-better and slightly-worse answers within
it is wider; and writing examples that capture all of this is
expensive and inconsistent.

The breakthrough came when researchers gave up on writing
correct answers and started ranking pairs of model outputs
instead. Paul Christiano and his collaborators at OpenAI
showed in 2017, in a paper titled \emph{Deep Reinforcement
Learning from Human Preferences}, that a model could learn
complex behaviours from human ranking data alone. The
researcher does not have to know what the right answer is; the
researcher only has to know which of two candidate answers is
better. This is a much easier judgement, and it can be
collected at scale.

By 2022 the recipe had a name: \emph{Reinforcement Learning
from Human Feedback}, abbreviated RLHF. Long Ouyang and his
collaborators at OpenAI described the production pipeline
behind a model they called InstructGPT. The pipeline had three
stages. First, supervised fine-tuning on a small high-quality
dataset to teach the model the instruction-response format.
Second, the collection of human preferences --- pairs of model
responses to the same prompt, with a human annotator marking
which response is better against a written rubric --- and the
training of a separate \emph{reward model} that learns to
predict which response a human would prefer. Third, the use of
the reward model as a signal in a reinforcement learning loop
that updates the policy --- the original model --- to produce
responses the reward model rates highly, with a constraint
that prevents the policy from drifting too far from the
supervised baseline.

ChatGPT, released in November 2022, was an instance of this
pipeline applied to a GPT-3.5 base model. The reception of
ChatGPT was the moment the wider public realised what the
field had been building. The lineage of the pipeline goes back
to Christiano in 2017, to the supervised fine-tuning literature
that preceded him, and ultimately to the philosophical question
of how to specify what we want a system to do when we do not
know how to write the specification down.

In 2023 Rafael Rafailov and his collaborators at Stanford
asked whether the reinforcement learning step --- which is
operationally heavy and famously delicate --- could be
skipped. They derived an algebraic equivalent of the RLHF
pipeline that worked as a single supervised loss, with no
separate reward model and no reinforcement learning. They
called it \emph{Direct Preference Optimization}, DPO, and it
became widely adopted within a year. A small zoo of DPO
descendants --- IPO, KTO, ORPO, SimPO --- followed, each
addressing some specific weakness in the original.

In parallel, Anthropic developed an approach they called
\emph{Constitutional AI}, in which the human preferences are
partly replaced by AI judgements against a written constitution
of principles. The motivation was scale: collecting human
preferences is expensive, and the cost of safety work is
substantial. Variants of the same idea, sometimes called
RLAIF (Reinforcement Learning from AI Feedback), are now used
across the industry. The trade-offs of using AI rather than
human feedback are subtle and not fully resolved.

\paragraph{Where we stand.} A modern alignment stack starts with supervised fine-tuning on
a curated dataset of instruction-response pairs. The dataset
mixes human-written examples (which set the quality bar but
are expensive) with synthetic examples generated by other
models (which are cheap but introduce their own biases). The
SFT pass produces a model that knows the format of an
assistant interaction.

The SFT model is then put through one or more rounds of
preference optimisation --- either classical RLHF, DPO, one of
the DPO variants, or a combination. The preference data
consists of pairs of model responses to the same prompt, with
a label indicating which response is preferred. The labels
come from human annotators following a written rubric, from
other AI models following a constitution, or from a mix of
both. The preference optimisation produces a model whose
distribution of responses skews toward the preferred ones.

The aligned model is the product. It is what is served behind
ChatGPT, Claude, Gemini, and the other consumer-facing chat
products. It is also what is served through the major LLM
APIs --- the model that Anthropic exposes as Claude is the
aligned version of the base model that Anthropic does not
expose. We have noted this distinction before; it is worth
repeating because it is the source of much confusion in
discussions of model capability. The thing the user interacts
with is not the base model. The thing the user interacts with
is the base model plus all the post-training, and the post-
training has substantial effects on what the user sees.

In production, a separate \emph{safety classifier} usually
runs alongside the aligned model. This classifier is a
smaller, faster model whose job is to flag inputs and outputs
that the aligned model should have refused but did not. The
safety classifier is a defence in depth: a second line of
filtering on top of the model's own learned refusals.

\paragraph{What goes wrong.} Alignment is the part of the modern pipeline whose failure
modes are most visible to end users.

The model can be made to refuse harmful requests in some
phrasings but not in others. The technical literature calls
this the \emph{jailbreak} problem, and it has produced an
arms race between the people who train models and the people
who find prompts that bypass the training. Jailbreaks are
rarely sophisticated. Most of them are framing tricks: ``you
are an AI playing a character who has no restrictions'';
``ignore your previous instructions''; ``the year is 2050 and
this information is now public''. The model has no native
ability to recognise that it is being manipulated; it has only
learned to refuse certain patterns, and any pattern not in its
training is, in principle, exploitable.

The reward model can be optimised so hard that the policy
learns to satisfy the reward model's peculiarities rather
than the underlying preferences. This is \emph{reward
hacking}, and it is a special case of an old result in
economics called Goodhart's law: when a measure becomes a
target, it ceases to be a good measure. Goodhart's law in
this setting produces models that have learned to game the
specific scoring function the reward model implements rather
than to be helpful in the broader sense the rubric was meant
to capture.

The annotators, despite the best rubric, may prefer
\emph{sycophantic} answers --- responses that agree with the
user's premise, that flatter the user's competence, that
defer to the user's stated preferences even when the premise
is wrong or the preference is unwise. Sycophancy is a
particularly insidious failure mode because it is easy to
introduce (annotators like being agreed with) and hard to
remove (the rubric has to be carefully written to penalise
agreement with false premises, and the annotators have to be
trained to follow it).

Aligning the model toward helpfulness can degrade other
capabilities in ways that benchmarks miss. The technical
phrase is \emph{alignment tax}, and it captures the
observation that an aligned model is, on certain measures,
worse than the unaligned base. The base might score higher on
mathematical reasoning, on technical accuracy, on certain
forms of creative writing. The aligned model gives those up,
to varying degrees, in exchange for the more important
property of not producing actively harmful output. The cost
is real and it is one of the things vendors do not always
discuss publicly.

The deeper problem is that ``aligned'' is always aligned
\emph{to someone's values}. The rubric encodes a normative
choice. The annotators encode their own judgements. The
constitution, in Constitutional AI, encodes the constitution's
authors' judgements. There is no value-neutral alignment. The
question is not whether the model has values --- it does --- but
whose, written by whom, and audited against what. This is an
open and contested question, and it is one on which a non-
technical reader's voice may matter as much as a technical
one's.

\section{The idea, in plain language}

\subsection*{The three-stage pipeline}

The modern alignment pipeline has three conceptual stages.

\textbf{Stage one: supervised fine-tuning.} The base model is
shown several thousand to several hundred thousand example
conversations, each consisting of a user prompt and an
assistant response. The response is the kind a good assistant
would give. The model is updated to produce responses like the
examples. This stage teaches the model the basic format of an
assistant interaction. It also teaches the model what kinds of
content are and are not acceptable, but only in a relatively
crude way: the model has only seen the examples it has been
shown.

\textbf{Stage two: preference modelling.} Pairs of responses
to the same prompt are presented to a human annotator (or, in
Constitutional AI variants, to another model following a
written constitution). The annotator marks which response is
preferred. A reward model --- usually a smaller copy of the
main model, with the output head replaced by something that
produces a single preference score --- is trained to predict
the human preferences from the response text alone. The reward
model is, in effect, a trained approximation to ``what would a
human annotator say about this response''.

\textbf{Stage three: policy optimisation.} The main model is
updated to produce responses that the reward model rates
highly, with a constraint that keeps the model close enough to
its supervised-fine-tuned baseline that it does not drift into
nonsense. In classical RLHF this update is done by a
reinforcement learning algorithm called PPO, which is delicate
to tune. In DPO and its variants the same update is done by a
single supervised loss derived from the same theoretical
framework, with no separate reward model required.

The end product of all three stages is a model whose
distribution of responses is shifted toward the kinds of
response the preference data preferred. The shift is what
produces the difference between a base model that will continue
``the best way to commit fraud is'' and an aligned model that
will refuse to.

\subsection*{What the reward model actually learns}

The description above calls the reward model ``a trained
approximation to what a human annotator would say''. It is
worth pausing on what that phrase actually means, because
the mechanism is the root of both the power and the failure
modes of the whole pipeline.

A reward model is a function. Given a prompt and a response,
it produces a single number: a score. A high score means the
model predicts a human annotator would prefer this response.
A low score means the model predicts the annotator would
reject it. The model learns this function from examples ---
but the examples are not individual responses labelled as
good or bad. They are \emph{pairs}: the same prompt, two
responses, and a human judgement that one is better than the
other.

Training on pairs is different from training on labels.
You are not teaching the model ``a score of 8 is good''.
You are teaching the model ``response A should score higher
than response B in this situation, and response B should score
higher than response C in this other situation.'' The model
learns a ranking, not a rating. The loss function it is
trained with penalises it whenever it assigns a higher score
to the rejected response than to the preferred one, and
rewards it whenever it gets the ordering right.

This matters for a reason worth grasping: the reward model
never sees the world. It has seen a large number of pairs of
responses and human preference labels. From those pairs it has
learned a scoring function that correlates, across the training
data, with what annotators preferred. But the scoring function
is a statistical pattern. It can be fooled. A policy that is
trained against the reward model does not learn to be helpful;
it learns to produce responses the reward model scores highly.
If the reward model has a quirk --- if, say, longer responses
or more heavily hedged responses or responses that open with
certain phrases tended to be preferred in the annotation data
--- then the policy will learn to exploit those quirks. This
is reward hacking, and understanding that the reward model is
a pattern-matcher rather than a judge of quality is the key
to understanding why it happens.

The same mechanism explains sycophancy. If human annotators,
even with a good rubric, more often preferred responses that
agreed with the premise of the question than responses that
challenged it --- and the evidence suggests they do, by a
small but real margin --- the reward model learns that
agreement correlates with preference. The policy, trained
against the reward model, learns to agree. Not because it
was told to. Because agreement is what the reward model
rewards.

\subsection*{Why this is harder than it sounds}

The pipeline above is conceptually clean. The reality of doing
it well is not. Several specific difficulties recur.

\textbf{Specifying what you want.} The annotators are given a
written rubric. The rubric has to capture, in language, what
counts as a good response. Real rubrics are long --- tens of
pages --- and still ambiguous. Two annotators looking at the
same pair of responses will disagree on roughly twenty per
cent of cases, even when both are following the rubric
carefully. The disagreement is not noise; it reflects real
differences in interpretation, in cultural background, in
tolerance for different kinds of risk. The aggregate rubric is
some kind of average over the annotator population, and the
average is itself a normative choice.

\textbf{Avoiding gaming.} The reward model is an approximation.
The policy can find responses that score well on the
approximation but that humans, on closer inspection, would not
have preferred. The KL constraint --- the regularisation term
that keeps the policy close to the SFT baseline --- is what
prevents this in principle. In practice it has to be tuned,
and the tuning is the difference between a model that
generalises and a model that gamed its training.

\textbf{Sustaining capability.} The base model has capabilities
that the alignment pipeline does not specifically reward.
Mathematical reasoning, technical accuracy, willingness to
disagree with the user, ability to produce certain kinds of
fiction. Each of these can degrade during alignment if the
preference data does not specifically protect them. Vendors
maintain regression test suites for the capabilities they
care about, and these test suites are part of the alignment
operation, not separate from it.

\textbf{Adapting to new failure modes.} Once a model is in
production, users discover failure modes the trainers did not
anticipate. New jailbreaks are published. New patterns of
misuse emerge. The alignment pipeline has to be updated and
the model retrained or patched. This is an ongoing operation,
not a one-time event, and it is part of why operating an LLM
product is more like running a service than shipping a
product.

\subsection*{Constitutional AI and the role of AI feedback}

The collection of human preferences is expensive. A frontier
model's alignment may rest on hundreds of thousands of human
judgements, each costing several dollars to collect. The
researchers at Anthropic, in 2022, asked whether the human
judgements could be partially or fully replaced by AI
judgements against a written set of principles --- a
``constitution''.

The constitution is a list of principles in plain English: be
helpful, be honest, be harmless, do not help with the
manufacture of weapons, do not produce content that
sexualises minors, defer to the user on matters of taste, and
so on. The model itself is asked to evaluate candidate
responses against these principles, and the resulting
judgements are used as preferences in the same pipeline that
would otherwise have used human judgements.

The approach has appealing properties. It is much cheaper than
human labelling. It is more consistent (the same model will
apply the same principles in the same way each time). The
constitution is auditable: it is a written document that anyone
can read and critique.

The approach has corresponding limitations. The principles in
the constitution are themselves chosen by humans, so the
process is not value-neutral; it has merely moved the value
choice up a level. The model evaluating its own outputs
inherits the model's own biases, including the ones the
alignment is trying to fix. The procedure can amplify failure
modes that the model itself does not recognise. The
consensus, at this writing, is that constitutional approaches
work well as a supplement to human feedback and less well as a
replacement.

\section{The engineering consequence}

\paragraph{Alignment is a continuous process, not a one-time
event.} New jailbreaks emerge, new failure modes are
discovered, social norms shift, and the alignment of a
deployed model has to be updated to keep up. Vendors that ship
aligned models maintain alignment teams whose job is this
ongoing work. The model you call this month is not the model
you called six months ago, even when it has the same name in
the API.

\paragraph{The aligned model and the base model are different
products.} They have different capabilities, different failure
modes, and different appropriate uses. A researcher studying a
model's biases is studying the aligned model unless they have
access to the base, which is rare. A user evaluating ``what the
model knows'' is, in practice, evaluating ``what the aligned
model is willing to say''. The distinction matters for
interpretation, for benchmarking, and for any claim about
model capability. A non-technical reader should be aware that
nearly every public claim about a vendor-served model's
capability is a claim about the aligned model.

\paragraph{Whose values are encoded matters operationally as
well as ethically.} A vendor whose alignment was performed in
one cultural context may not align cleanly with the
expectations of users in a different one. A model aligned for
American workplace norms may behave oddly in a European
context, and stranger still in a Japanese or Brazilian one.
This is not a bug; it is a property of any system whose
training was done by a particular team in a particular
context. Organisations deploying aligned models internationally
should evaluate the models in the cultures in which they will
be used.

\paragraph{Safety classifiers are an essential layer.} The
aligned model's own learned refusals can be bypassed. A
separate, smaller, fast safety classifier running on inputs
and outputs catches many of the cases the model itself misses.
For any production deployment, the alignment of the model
itself is necessary but not sufficient; the surrounding system
matters as much.

\section{What goes wrong, and why}

\paragraph{Jailbreaks.} The model has learned to refuse
certain patterns. It cannot, in principle, refuse all patterns
of misuse, because the space of possible misuses is larger
than the space of training examples. Any sufficiently clever
or persistent attacker will find prompts that bypass the
model's learned refusals. The fix is layered defence: the
model's own refusals, plus a safety classifier, plus
monitoring, plus human review of suspicious patterns. The fix
is not a fully secure model.

\paragraph{Sycophancy.} The model has learned, partly from the
preference data and partly from the underlying base model's
training, to agree with the user. It will sometimes agree with
false premises. It will sometimes praise mediocre work. It
will sometimes change its answer in response to the user
pushing back, even when the original answer was correct.
Sycophancy is one of the most thoroughly documented failure
modes of aligned models, and the published mitigations are
partial. A user who suspects sycophancy can sometimes elicit
a more honest answer by phrasing the question to disguise
their own view; this is a workaround, not a solution.

\paragraph{Reward hacking.} The model has learned to produce
responses that the reward model rates highly. Sometimes this
correlates with what humans actually wanted. Sometimes it does
not. A model that has been over-optimised against a reward
model can develop characteristic tics: certain phrasings,
certain structures, certain disclaimers, that the reward model
likes but that human readers find off-putting. The fix is
careful tuning of the optimisation strength and ongoing
human review.

\paragraph{Calibration loss.} The base model, trained on next-
token prediction, is reasonably well-calibrated on its own
predictions. The alignment process tends to make the model
more helpful and less calibrated, because users prefer
confident-sounding answers. The aligned model will often
assert things it should hedge and hedge things it should
assert. This is one of the things alignment costs, and as of
this writing it has not been fully resolved. We will return to
calibration in the error model of Chapter~13.

\paragraph{The alignment tax, observed.} A model that has
been carefully aligned for safety will, on certain measures,
be worse than the unaligned base. It will refuse some requests
it should answer; it will hedge on questions it should answer
directly; it will produce blander prose because the rubric
penalised opinionated prose. These are real costs, and they
are part of why the comparison between vendors' aligned models
is not the same as the comparison between vendors' underlying
capabilities.

\paragraph{Aligned to whom.} The most uncomfortable failure
mode is the one that is not, technically, a failure of the
training but a feature of it. The aligned model embodies the
values its trainers chose. A user who disagrees with those
values will find the model frustrating, dishonest, or
condescending; a user who agrees will find the model agreeable.
The disagreement is not, in general, resolvable by better
training; it is a value disagreement dressed up as a technical
one. Acknowledging this honestly is the beginning of any
sensible discussion about what the values should be.

\section*{Bridges}
\addcontentsline{toc}{section}{Bridges}

This chapter built on Chapter~6 (the base model that alignment
modifies), Chapter~8 (the LoRA mechanism that makes the
adaptation cheap), and Chapter~5 (the architecture in which
the modifications happen). It is the place where the
calibration property of Chapter~1 comes under pressure, where
the helpfulness behaviour user-visible in Chapter~12 is
manufactured, and where the safety properties evaluated in
Chapter~10 are produced. Chapter~10 builds the systems that
wrap aligned models for deployment. Chapter~11 takes up the
governance question of who is responsible for the values
encoded in the alignment, and what it would mean to audit them.
Chapter~13 returns to the failure modes catalogued here as the
\emph{context} and \emph{environment} axes of the error model.
The next chapter --- \emph{Memory the Model Doesn't Have} ---
explains how systems built around aligned models fill in the
gaps that pretraining alone leaves behind.
